%%%%%%%%%%%%%%%%%%%%%%%%%%%%%%%%%%%%%%%%%%%%%%%%%%%%%%%%%%%%%%%%%%%%%%
%  第 1 个逻辑月 · 月度报告
%  主题：小块：从抽象到具体（复现学长 + UDF 基线 + 算法提要）
%
%  说明：本报告由三人的既有材料按组长提纲的逻辑月重排而成。
%        尚无素材之处按规范留空并标注，不编造。
%  编译：xelatex 本文件（跑两遍）
%%%%%%%%%%%%%%%%%%%%%%%%%%%%%%%%%%%%%%%%%%%%%%%%%%%%%%%%%%%%%%%%%%%%%%
\documentclass[stage1,fontset=fandol]{ydzjreport}

\stageonesetup
  {第 1 个逻辑月 · 基础建立}%
  {冯旭东、陈思文、张昊洋}%
  {复合微通道分区冷却}%
  {【待确认】逻辑月 1 对应的自然月}%
  {小块划分；学长复现；水物性 UDF；算法提要}%
  {确立小块的划分方式，复现学长的建模与物理性质（UDF）思路，
   并完成控制方向的算法提要。
   }%
  
\begin{document}
\makestageonehead

% ==================================================================
\section{本月研究任务}

\subsection{任务目标}

按组长提纲，本月对应关键词为
\textbf{「复现论文+课题组经验的结构、了解物理性质，编写UDF，初步搭建冷板思路，初步构建控制流程」}、
据此分解为四项任务：

\begin{enumerate}
  \item 通过阅读文献确立\textbf{冷板的优化方式}，把空间不均匀热需求落实为研究的结构；
  \item 复现已有文献、学长在\textbf{结构建模}上的思路；
  \item 复现已有文献、学长在\textbf{物理性质（UDF）}上的设置，并核查其正确性；
  \item 初步控制方向完成\textbf{流程框架构建}。
\end{enumerate}

\subsection{完成情况}

\begin{table}[H]
\centering
\caption{本月任务完成情况}
\label{tab:tasks}
\begin{tabularx}{\textwidth}{@{}Y c Y@{}}
\toprule
\rowcolor{ydzjlight} 任务 & 状态 & 说明 \\
\midrule
小块划分方式确立 & 部分完成 & 已形成概念表述，划分图待补 \\
学长结构建模复现 & 完成 & 已复现论文中的三种流道布局 \\
学长物性（UDF）思路复现 & 已完成 & 核查后发现原函数存在误差 \\
控制算法提要 & 已完成 & 文献要点与二阶对象特性已整理 \\
\bottomrule
\end{tabularx}
\end{table}


% ==================================================================
\section{核心概念梳理}

\subsection{概念一：非均匀热载荷与小块划分}

芯片热流密度在空间上高度不均匀：结区集中发热，背景区几乎不发热。
若采用均匀结构，背景区会被「一视同仁」地强力冷却，形成\textbf{过度冷却}——
既拉大温度梯度，又浪费泵功。

\textbf{小块}的作用，是把连续冷板上的空间不均匀热需求，转化为
\textbf{可单独研究、可单独供液}的流道单元。

需要说明的是：各小块在连续冷板上形成不同几何区域，
\textbf{并不必然意味着把冷板机械切开或完全隔热}。
先研究单元、再研究组合，有利于识别各形态的作用，
并为独立供液或流量调控预留条件。

\textcolor{red}{【待补：小块划分图】
几何结构图}

\subsection{概念二：物理性质是计算的基线}

冷却工质水的物性随温度变化。若物性函数有误，
\textbf{所有温度与压差结果都不可信}，后续结构比较也失去意义。
因此\textbf{物性核查必须放在结构比较之前}。

\textcolor{red}{补充初步控制框架}

% ==================================================================
\section{阶段一：通读与摘录}

\subsection{学习路径}

\begin{enumerate}
  \item 先通读学长攻略（课题组已有经验）、结构类论文，梳理几何布局、进出口组织与边界条件；
  \item 复现其中的代表性流道，取得可对照的性能趋势；
  \item 转向物理性质：核查计算所用的水物性函数，确认随温度变化的正确性；
  \item 控制方向并线推进：整理控制文献要点与二阶对象的动态特性，形成算法提要。
\end{enumerate}

\subsection{关键摘录}

\begin{table}[H]
\centering
\caption{文献要点与本阶段工作}
\label{tab:notes}
\begin{tabularx}{\textwidth}{@{}c Y Y@{}}
\toprule
\rowcolor{ydzjlight} 文献 & 论文要点 & 本阶段的应用 \\
\midrule
Ao 等 & 方形、圆形与六边形微通道的比较 & 复现三种流道，学习短流路与局部强化的组织方式 \\
Yin 等 & 用平均努塞尔数、最高温度、压降与综合指标比较结构 & 借鉴其数据分析口径，为后续统一指标做准备 \\
Seong 等 & 倾斜肋与交错排列影响流动重分配 & 拓展倾斜、交错与复合结构的候选思路 \\

\bottomrule
\textcolor{red}{补充文献数据，扩大矩阵}
\end{tabularx}
\end{table}

% ==================================================================
\section{阶段二：导出与复现}

\subsection{复现目标}

本月复现两项内容：
\textbf{(1)} 学长论文中的方形、圆形与六边形三种微通道的结构与流动换热组织；
\textbf{(2)} 计算所用的水物性函数（UDF），核查其随温度变化的正确性。

\subsection{复现环境与步骤}

\textbf{环境}：ANSYS Fluent，水为工质，固体域与流体域耦合求解；
物性核查采用 Python 后处理。

\textbf{步骤}：

\begin{enumerate}
  \item 按论文的方形、圆形与六边形布局分别构造流道，梳理流体域、固体域与加热区域；
  \item 对照论文的边界条件与求解方法开展计算，检查残差与监测量；
  \item 对三种结构做趋势对照，关注热负荷与流量变化下的温度响应；
  \item 核查原水物性 UDF：在 298.15 K 校验点比对参考物性值；
  \item 记录误差量级，判断是否需要重新拟合。
\end{enumerate}

\subsection{复现结果}

\textbf{结构复现}：三种流道的复现已完成，计算结果与论文的趋势一致——
流量增加时结温总体下降、压降增加；热流密度增加时结温上升。

\begin{table}[H]
\centering
\caption{三种流道的复现情况}
\label{tab:repro}
\begin{tabularx}{\textwidth}{@{}Y Y Y@{}}
\toprule
\rowcolor{ydzjlight} 流道形态 & 复现关注内容 & 结果 \\
\midrule
方形微通道 & 方形流路及中心扰动结构的影响 & 已完成，趋势与论文一致 \\
圆形微通道 & 环形流路及流体分配的影响 & 已完成，趋势与论文一致 \\
六边形微通道 & 六边形布局及局部流动组织 & 已完成，趋势与论文一致 \\
\bottomrule
\end{tabularx}
\end{table}

【缺数据：逐点复现误差与网格敏感性统计】

\textbf{物性核查}：在 298.15 K 校验点，原水物性函数相对参考值
\textbf{密度偏低约 60.8\%、比热偏低约 10.5\%、黏度偏低约 26.7\%}。

这是一个\textbf{必须记录的重要发现}：它意味着此前所有基于该物性的
温度与压差结果都存在系统偏差，结构之间的排序也可能受影响。
本月据此完成了新物性 UDF 的编写与编译，挂接与复算列入后续月份。

【缺数据：新 UDF 挂接后的物性自检结果】

% ==================================================================
\section{阶段性收获}

\subsection{认知上的变化}

\textbf{一、从「整体冷却」转向「按需分区」。}
原来倾向于把冷板当作一个整体来强化，认为加密结构总归有利。
现在认识到：热流密度本身不均匀，均匀强化会在低负荷区造成过度冷却，
\textbf{冷却资源必须按区域配置}。

\textbf{二、物性不是背景条件，而是结论的前提。}
复核原物性函数之前，我们默认它是正确的。
发现偏差后才知道：\textbf{在物性不可信的情况下比较结构，等于在错的基准上排序}。
因此把「物理性质」放在结构比较之前的做法是必要的。

\subsection{尚存的疑问}

\begin{enumerate}
  \item 【待确认9】小块究竟按什么划分？三份材料里出现过
        「底部中央浮动热流区」「低温区／变化区／高温区」两种说法，是否为同一件事？
  \item 【待确认1】需要复现的学长具体是谁？我们复现的是已发表论文，
        与「学长」是否等同？
  \item 修正物性之后，此前三种流道的趋势是否仍然成立？
        需要复算才能判断，不能只重算后处理指标。
\end{enumerate}

% ==================================================================
\section{下月计划}

\begin{table}[H]
\centering
\caption{下月工作计划}
\label{tab:plan}
\begin{tabularx}{\textwidth}{@{}Y Y c@{}}
\toprule
\rowcolor{ydzjlight} 事项 & 预期产出 & 时间 \\
\midrule
挂接新水物性 UDF 并完成物性自检 & 物性校验记录 & 第 1 周 \\
补齐小块划分图，明确划分依据与数量 & 小块划分方案 & 第 1--2 周 \\
建立统一的评价指标体系 & 指标定义表 & 第 2--3 周 \\
小块内的形态与扰流单元仿真 & 形态对比数据 & 第 3--4 周 \\
\textcolor{red}{补充自控学习}\\
\end{tabularx}
\end{table}

% ==================================================================
\begin{thebibliography}{9}
\bibitem{ao} Ao C, Xu B, Chen Z. Design and performance optimization of novel microchannel heat sinks for high-power electronic chips[J]. International Communications in Heat and Mass Transfer, 2026, 178: 111803.
\bibitem{yin} Yin Y, Wang Z, Ma S, et al. A novel coaxial open annular TSV pin-fin embedded in microchannels[J]. International Communications in Heat and Mass Transfer, 2026, 172: 110619.
\bibitem{seong} Seong H B, Gwon J G, Lee M S, et al. Numerical study and ANN analysis on heat transfer and flow characteristics of microchannel heat sink with inclined rib--pin fin[J]. International Communications in Heat and Mass Transfer, 2026, 178: 111593.
\bibitem{yuan} 袁永帅, 等. 当前芯片散热技术方法和特点研讨[J]. 2024.
\bibitem{xiong} 熊园园, 等. 基于内嵌式微通道芯片散热结构设计研究综述[J]. 2024.
\end{thebibliography}

\end{document}
